% Chapter 5: BCFW Recursion
% Based on research from 2026-02-17

\chapter{On-Shell Recursion (BCFW)}
\label{ch:bcfw}

\section{The BCFW Shift}

Choose two legs $i$ and $j$. Deform their spinors:
\begin{align}
\lambda_i &\to \lambda_i + z \lambda_j \\
\tilde{\lambda}_j &\to \tilde{\lambda}_j - z \tilde{\lambda}_i
\end{align}

Under this shift:
\begin{equation}
p_i(z) + p_j(z) = p_i + p_j \quad \text{(momentum conserved)}
\end{equation}

\section{The Recursion Formula}

If $A_n(z) \to 0$ as $z \to \infty$, then by Cauchy's theorem:
\begin{equation}
\boxed{A_n = \sum_{\text{channels } P} A_L^{\text{on-shell}} \frac{1}{P^2} A_R^{\text{on-shell}}}
\end{equation}

where the sum runs over all factorization channels where $P$ is on-shell ($P^2 = 0$).

\section{Why It Works}

Traditional Feynman diagrams explode combinatorially:
\begin{center}
\begin{tabular}{|c|c|}
\hline
$n$ gluons & Feynman diagrams \\
\hline
4 & 4 \\
5 & 25 \\
6 & 220 \\
10 & $\sim 10^7$ \\
\hline
\end{tabular}
\end{center}

BCFW recursion reduces this to $O(n^3)$ complexity.

\section{Application to Single-Minus}

The arXiv:2602.12176 paper uses BCFW with:
\begin{itemize}
\item All-plus reference leg
\item Single-minus shifted leg
\item Induction on $n$
\end{itemize}

This proves the nonvanishing of single-minus amplitudes for all $n$.

\section{Soft Limits}

As $p_n \to 0$:
\begin{equation}
A_n \to \frac{\langle n-1, n \rangle}{\langle n, 1 \rangle} A_{n-1}
\end{equation}

This provides a consistency check on amplitude formulas.
